\documentclass[11pt,a4paper]{article}
\usepackage[T1]{fontenc}
\usepackage{lmodern,amsmath,amssymb,amsthm,booktabs,microtype,enumitem}
\usepackage[a4paper,margin=26mm]{geometry}
\usepackage{xcolor}
\definecolor{ink}{HTML}{183747}
\definecolor{plum}{HTML}{6A3151}
\usepackage[colorlinks=true,allcolors=plum]{hyperref}
\hypersetup{pdftitle={When a more cyclic majority tournament is more probable},
 pdfauthor={Anonymous},
 pdfsubject={Random elections, majority tournaments, paired comparisons}}
\newtheorem{theorem}{Theorem}[section]
\newtheorem{lemma}[theorem]{Lemma}
\newtheorem{proposition}[theorem]{Proposition}
\newtheorem{corollary}[theorem]{Corollary}
\theoremstyle{definition}
\newtheorem{remark}[theorem]{Remark}
\newtheorem{definition}[theorem]{Definition}
\DeclareMathOperator{\lin}{lin}
\newcommand{\E}{\mathbb E}
\newcommand{\R}{\mathbb R}
\newcommand{\dd}{\mathrm d}
\title{\color{ink}\bfseries When a more cyclic majority tournament is more probable}
\author{Anonymous\\[3pt]{\small DOI \href{https://doi.org/10.5281/zenodo.24526810}{10.5281/zenodo.24526810}}}
\date{November 2026}
\begin{document}
\maketitle

\begin{abstract}
When many voters rank a fixed set of candidates independently and uniformly at random (``impartial culture''), the head-to-head majorities form a random tournament. Harrison-Trainor conjectured that, in the limit of many voters, a labelled tournament with fewer cyclic triples (equivalently, a larger sum of squared win counts) is always the more probable one. We show that the conjecture is false. For every odd $t$ there are tournaments on $14t$ candidates in which the tournament with $3t^3$ more cyclic triples is $2^{3t^2+o(t^2)}$ times as probable. The reason is that the probability cost of a tournament is governed by pairs of candidates whose strengths must be nearly equal, while the cyclic-triple count is governed by triples, and the two counts can order the same block structures oppositely. A rigorous rational certificate gives an explicit pair on 78 candidates: the tournament with three more cyclic triples is more than $10^5$ times as probable, and the reversal persists for the probabilities of isomorphism classes. Monte Carlo computation, which is not a proof, finds no reversal on seven or fewer candidates and finds reversals on 13 and 14 candidates. The extremal half of the conjecture is true: a compression argument shows that the transitive tournaments are exactly the most probable ones, for every number of candidates and for a wide class of paired-comparison models. A second part asks how much the win totals of a round robin reveal. For any fixed latent-type comparison model with probabilities strictly between zero and one, either the win totals determine every outcome's probability (the Bradley--Terry case), or a typical outcome shares its win totals with outcomes exponentially more and exponentially less probable on the scale of the number of pairs. A conditional test built from directed triangles is exactly calibrated at every vector of win totals and detects every fixed departure of this kind from a single round robin. Finally, one observation per pair cannot support the minimax Bradley--Terry--Luce test of Makur and Singh.
\end{abstract}

\section{Introduction}\label{sec:intro}

This section states the problem, explains why it matters, summarises what was known and what is new, and gives a roadmap.

\subsection{The question}

Suppose $N$ voters each rank $n$ labelled candidates, every one of the $n!$ rankings being equally likely and the voters independent. This \emph{impartial culture} is the standard null model of social choice: it is the model under which classical paradoxes of voting are counted and against which voting rules are stress-tested. For each pair of candidates, the majority prefers one to the other when $N$ is odd, so the election produces a \emph{tournament}: a complete directed graph on the candidates in which $i\to j$ means that a majority ranks $i$ above $j$. Condorcet's paradox is the observation that this tournament can contain a directed triangle $i\to j\to k\to i$, a \emph{cyclic triple}.

For a prescribed labelled tournament $T$ let $p(T)$ be the limit, as $N\to\infty$ through odd integers and with $n$ fixed, of the probability that the majority tournament is exactly $T$. Write $d_v$ for the number of wins (outdegree) of candidate $v$. Harrison-Trainor \cite{harrison} called $\lin(T)=\sum_v d_v^2$ the \emph{linearity} of $T$ and proposed the following.

\begin{quote}
\textbf{Conjecture} (Harrison-Trainor; Conjecture~6.3 of the arXiv version of \cite{harrison}). For labelled tournaments $T,U$ on the same candidates, $\lin(T)>\lin(U)$ implies $p(T)>p(U)$.
\end{quote}

Linearity is a disguised count of cyclic triples. Every transitive triple has exactly one vertex beating the other two, so the number $c(T)$ of cyclic triples satisfies
\begin{equation}
 c(T)=\binom{n}{3}-\sum_v\binom{d_v}{2}
 =\binom{n}{3}-\tfrac12\lin(T)+\tfrac12\binom{n}{2}.
 \label{eq:triangles}
\end{equation}
This is the classical count of cyclic (circular) triads on which Kendall and Babington Smith's test for paired comparisons is based (see \cite[Section~4.1.1]{SD2026} for a current account). The conjecture therefore says: \emph{the fewer cyclic triples a majority tournament has, the more likely it is.} It is a natural strengthening of the intuition that impartial culture favours orderly outcomes.

\subsection{Why it matters and what was known}

The distribution of the majority tournament under impartial culture is the reference distribution for many questions in social choice, such as the probability of a Condorcet winner, of a transitive outcome or of a given tournament solution, and it is the distribution sampled when majority graphs are generated for simulations. Knowing which individual outcomes are likely, and why, describes the shape of this distribution beyond its classical summary statistics.

As $N\to\infty$, the normalised pairwise margins converge to a Gaussian vector in which two comparisons sharing a candidate have covariance $\pm1/3$ and comparisons of disjoint pairs are uncorrelated; see \cite{harrison} and \cite{close}. Harrison-Trainor used this limit to compute $p(T)$ exactly for three and four candidates and numerically for five, and verified the conjecture in those cases \cite{harrison}. Mossel showed that the probability that the majority tournament is transitive is $\exp(-\Theta(n^{5/3}))$ \cite[eq.~(6)]{Mossel2010}. The limiting probabilities coincide with the marginal likelihood of a Thurstone--Mosteller paired-comparison model with Gaussian random strengths, a model studied by Varin and Firth \cite{varin}. Close elections, in which all margins are conditioned to be small, behave differently: there every tournament is asymptotically equally likely \cite{close}. We study the unconditional probability of the unweighted majority graph.

\subsection{What is new}

\begin{enumerate}[leftmargin=*,label=(\alph*)]
\item \textbf{The conjecture is false} (Theorem~\ref{thm:headline}). The pairs are ``ordered block'' tournaments: candidates are split into blocks, every member of an earlier block beats every member of a later block, and each block is internally regular (every member wins half its internal games). Splitting $14t$ candidates as $(7t,7t)$ or as $(t,9t,3t,t)$ gives the more cyclic second tournament a probability larger by a factor $2^{3t^2+O(t\log t)}$.
\item \textbf{The mechanism.} To realise a block structure, the latent strengths of the candidates must separate the blocks, which costs only $O(n\log n)$ in log probability, while every pair inside a block costs a factor of about two. So log probability is governed by $\sum_j b_j^2$, the number of internal \emph{pairs}, whereas cyclic triples lie inside blocks and are governed by $\sum_j b_j^3$, the number of internal \emph{triples} (Proposition~\ref{thm:blocks}). These two sums can rank partitions oppositely.
\item \textbf{An explicit certified counterexample} on 78 candidates (Proposition~\ref{prop:78}) with a factor above $10^5$, which also reverses the probabilities of the two isomorphism classes (Corollary~\ref{cor:classes}).
\item \textbf{Small cases} (Section~\ref{sec:small}). Monte Carlo integration, which is numerical evidence and not a proof, finds no reversal among all isomorphism types on three to seven candidates (456 types on seven), finds no reversal among block tournaments on up to 12 candidates, and finds reversals on 13 candidates (by about 1\%) and on 14 candidates (by a factor of about 2.7), the latter including the $t=1$ member of the family above.
\item \textbf{The extremal case is true} (Theorem~\ref{thm:compression}). The transitive tournaments are exactly the most probable labelled tournaments, for every $n$ and for every exchangeable latent-strength model whose win probabilities increase strictly with strength.
\item \textbf{What win totals reveal} (Part~II). The conjecture is a statement about how well the vector of win totals predicts probabilities. We show that for any fixed latent-type comparison model, either the win totals determine probabilities exactly (Bradley--Terry case; qualitatively due to Lauritzen and to Iyer) or typical outcomes have same-score neighbours $e^{\pm cn^2}$ times as probable (Theorem~\ref{thm:general-kernels}). An exactly calibrated triangle test detects every such model from one round robin (Section~\ref{sec:tests}), and one observation per pair cannot support minimax Bradley--Terry--Luce testing (Proposition~\ref{thm:one-rep}).
\end{enumerate}

\paragraph{What is not claimed.} All election results concern the limit of many voters at a fixed number of candidates; we give no electorate size at which a given finite reversal appears. The smallest number of candidates admitting a reversal is not determined: the certified bound is 78, and the numerical evidence places it between 8 and 13. The statistical results of Part~II are pointwise in a fixed model; apart from Proposition~\ref{thm:one-rep} they give no uniform detection rate.

\subsection{Roadmap}

Part~I concerns elections. Section~\ref{sec:gauss} recalls the Gaussian representation of $p(T)$. Section~\ref{sec:mode} proves a two-sided bound on $\log p(T)$ valid for every tournament. Section~\ref{sec:blocks} computes both statistics for ordered block tournaments and proves Theorem~\ref{thm:headline}. Section~\ref{sec:finite} gives the certified 78-candidate example and the small-case computations. Section~\ref{sec:compression} proves that transitive tournaments maximise probability. Part~II (Sections~\ref{sec:general-kernels}--\ref{sec:onerep}) treats general comparison models, conditional tests and replication. Section~\ref{sec:discussion} discusses scope, uses and open problems. Appendix~\ref{app:certificate} records the rational arithmetic of the certificate and Appendix~\ref{app:code} the companion programs.

\paragraph{Notation.} Logarithms are natural. $\Phi$ and $\phi$ are the standard normal distribution function and density, $\ell=\log\Phi$, and $r=\ell'=\phi/\Phi$. A tournament on vertex set $[n]=\{1,\dots,n\}$ is written as its set of arcs $i\to j$; $d_v$ is the outdegree of $v$; $\binom{n}{2}$ is the number of pairs.

\part*{Part I. Majority tournaments under impartial culture}
\addcontentsline{toc}{part}{Part I}

\section{A Gaussian model for the majority margins}\label{sec:gauss}

This section rewrites $p(T)$ as the expectation of a product of normal probabilities, the form used in every later proof.

Let $X_1,\ldots,X_n$ and $Z_{ij}$ for $i<j$ be independent standard normal variables, and put $Z_{ji}=-Z_{ij}$. Then the limiting majority tournament has the same law as the tournament defined by
\begin{equation}
 i\longrightarrow j\quad\Longleftrightarrow\quad X_i-X_j+Z_{ij}>0.
 \label{eq:representation}
\end{equation}
Indeed, encode one ballot by signs $Y_{ij}=1$ if it ranks $i$ above $j$ and $-1$ otherwise. These signs have mean zero and variance one; disjoint comparisons are uncorrelated; and among the six orders of three candidates, $Y_{ij}$ and $Y_{ik}$ agree four times, giving covariance $1/3$ (reversing one comparison changes the sign). The vector with coordinates $(X_i-X_j+Z_{ij})/\sqrt{3}$ has exactly these covariances. By the fixed-dimensional central limit theorem the normalised margins converge to this Gaussian vector, which gives zero probability to a tied margin; this proves~\eqref{eq:representation}. The same covariance structure is the starting point of \cite{harrison,close}.

Conditioning on $X=(X_i)_i$ makes the comparisons independent, with $\Pr(i\to j\mid X)=\Phi(X_i-X_j)$. Hence
\begin{equation}
 L_T(x)=\sum_{i\to j}\ell(x_i-x_j),\qquad
 p(T)=\E\exp L_T(X).
 \label{eq:conditional}
\end{equation}
Thus $p(T)$ is the Gaussian-prior marginal likelihood of the Thurstone--Mosteller (probit) paired-comparison model for one complete round robin; Varin and Firth \cite[eqs.~(3)--(5)]{varin} study this likelihood for estimation. Here it is used to compare probabilities of entire tournaments. We call $X_i$ the \emph{strength} of candidate $i$.

\section{A two-sided bound for every tournament}\label{sec:mode}

This section proves that $\log p(T)$ equals the maximum of a concave function up to an error of order $n\log n$, uniformly over all tournaments. Define the penalised log likelihood and its maximum,
\[
 H_T(x)=L_T(x)-\tfrac12\|x\|^2,\qquad M_T=\max_{x\in\R^n}H_T(x).
\]

\begin{theorem}\label{thm:mode}
For every tournament $T$ on $n$ vertices,
\begin{equation}
 M_T-\tfrac{n-1}{2}\log(n+1)\ \le\ \log p(T)\ \le\ M_T .
 \label{eq:mode}
\end{equation}
In the lower bound $M_T$ may be replaced by $H_T(x)$ at any vector $x$.
\end{theorem}

\begin{proof}
Let $Z$ be standard normal. Completing the square gives $\log\E[e^{uZ}\mid Z\le t]=u^2/2+\ell(t-u)-\ell(t)$, and differentiating twice in $u$ at $u=0$ gives $1+\ell''(t)=\operatorname{Var}(Z\mid Z\le t)>0$. Also $\ell''(t)=-r(t)(t+r(t))<0$, since the conditional mean $-r(t)$ of $Z$ given $Z\le t$ is below $t$. Hence $-1<\ell''<0$, and the Hessian of $-H_T$ lies, as a quadratic form, between $I$ and $I+K$, where
\[
 K=\sum_{i<j}(e_i-e_j)(e_i-e_j)^{\mathsf T}=nI-\mathbf1\mathbf1^{\mathsf T}
\]
is the Laplacian of the complete graph. As $H_T(x)\le-\|x\|^2/2$, the maximiser exists and is unique. Taylor's formula at the maximiser bounds $H_T$ above and below by quadratics with common maximum $M_T$ and precision matrices $I$ and $I+K$. Integrating the two Gaussian functions with the factor $(2\pi)^{-n/2}$ from~\eqref{eq:conditional}, the upper one contributes $1$ and the lower one $\det(I+K)^{-1/2}=(n+1)^{-(n-1)/2}$. For the last sentence, apply the same lower Taylor bound at an arbitrary point $x$: the extra linear term integrates to a factor at least one by Jensen's inequality.
\end{proof}

The error term cannot be reduced in general, as regular tournaments show. A tournament on odd $n$ is \emph{regular} if every vertex has outdegree $(n-1)/2$.

\begin{proposition}\label{prop:regular}
Uniformly over regular tournaments on odd $n$,
\[
 \log p(T)=-\binom n2\log 2-\tfrac{n-1}{2}\log n+O(n).
\]
\end{proposition}

\begin{proof}
Regularity gives $\sum_{i\to j}(x_i-x_j)=0$, so concavity of $\ell$ gives $M_T=L_T(0)=-\binom n2\log 2$, and the lower estimate follows from~\eqref{eq:mode}. For the upper estimate put $D(z)=\ell(0)+\ell'(0)z-\ell(z)$ and $g(z)=\min\{z^2,|z|\}$. There is an absolute constant $c>0$ with $D\ge cg$: indeed $D(0)=D'(0)=0$, $D''=-\ell''$ has a positive minimum on $[-1,1]$, and convexity beyond the endpoints gives $D(z)\ge |z|\min\{D(-1),D(1)\}$ for $|z|\ge1$. Since the linear terms cancel, $L_T(x)\le M_T-c\sum_{i<j}g(x_i-x_j)$. By the arithmetic--geometric mean inequality,
\[
 e^{-c\sum_{i<j}g(x_i-x_j)}
 \le \frac1n\sum_{j=1}^n\prod_{i\ne j}e^{-(cn/2)g(x_i-x_j)}.
\]
For standard normal $X$, uniformly in $y$, $\E e^{-(cn/2)g(X-y)}\le C/\sqrt n$: bound the normal density by its maximum and split the integral at $|X-y|=1$. Conditional on $X_j$, the $n-1$ factors in the $j$-th term are independent, so $p(T)\le e^{M_T}(C/\sqrt n)^{n-1}$.
\end{proof}

For comparison, Mossel's estimate \cite{Mossel2010} says that the $n!$ transitive labellings together have probability $\exp(-\Theta(n^{5/3}))$, so each has $\log p(T)=-\Theta(n^{5/3})-\log n!$, whereas a regular tournament has $\log p(T)\approx-(\log 2)n^2/2$.

\section{Ordered regular blocks and the disproof}\label{sec:blocks}

This section computes the two competing statistics for block-structured tournaments and deduces the counterexample family.

\begin{definition}
Let $b_1,\ldots,b_k$ be odd positive integers with sum $n$. An \emph{ordered regular block tournament} with block sizes $(b_1,\dots,b_k)$ partitions $[n]$ into consecutive blocks of these sizes, listed from strongest to weakest; every vertex of an earlier block beats every vertex of a later block, and each block carries a regular tournament. One explicit choice on a block labelled by $\mathbb Z_b$ is the \emph{circulant} tournament with $u\to v$ when $v-u\bmod b\in\{1,\ldots,(b-1)/2\}$. Put $m=\sum_j\binom{b_j}{2}$, the number of pairs inside blocks.
\end{definition}

\begin{proposition}\label{thm:blocks}
For fixed $k$, uniformly over all such block sizes and internal regular tournaments,
\begin{align}
 \log p(T)&=-m\log 2+O_k(n\log(n+1)),\label{eq:blockprob}\\
 \lin(T)&=\sum_{r=0}^{n-1}r^2-\frac1{12}\sum_j(b_j^3-b_j).\label{eq:blocklin}
\end{align}
\end{proposition}

\begin{proof}
Inside a regular block, $\sum_{i\to j}(x_i-x_j)=0$ because each vertex occurs equally often with each sign, so concavity of $\ell$ bounds the internal log probabilities of block $j$ by $\binom{b_j}{2}\ell(0)$. Cross-block terms are nonpositive. Hence $L_T\le -m\log 2$ and $p(T)\le 2^{-m}$.

For the lower bound put $x_i=(\frac{k+1}{2}-j)h$ on block $j$, with $h=\sqrt{2\log(n+1)}$. Internal differences vanish and cross-block differences are at least $h$. The Gaussian tail bound $1-\Phi(h)\le e^{-h^2/2}=1/(n+1)$ gives $\ell(h)\ge -\log(1+1/n)\ge-1/n$, so $L_T(x)\ge-m\log 2-n/2$, while $\|x\|^2=O_k(n\log(n+1))$. The lower bound of Theorem~\ref{thm:mode} proves~\eqref{eq:blockprob}.

For~\eqref{eq:blocklin}, let $s$ be the number of vertices in the blocks after (weaker than) a block of size $b$. All vertices of that block have score $s+(b-1)/2$. Ordering the block transitively instead would give them the scores $s,\ldots,s+b-1$, and the sum of squares would increase by the sum of squared deviations of $0,\ldots,b-1$ from their mean, namely $b(b^2-1)/12$. Ordering every block transitively gives each score $0,\dots,n-1$ once.
\end{proof}

If $b_j/n\to a_j$, the two statistics therefore behave as
\[
 \frac{\log p(T)}{n^2}\longrightarrow
 -\frac{\log2}{2}\sum_j a_j^2,\qquad
 \frac{\lin(T)}{n^3}\longrightarrow
 \frac13-\frac1{12}\sum_j a_j^3 .
\]
Pairs inside a block cost probability even though their endpoints may have equal strengths, while cyclic triples lie entirely inside blocks. A partition with a smaller sum of squares but a larger sum of cubes is more probable and more cyclic.

\begin{theorem}\label{thm:headline}
For every positive odd $t$ let $A_t$ and $B_t$ be ordered regular block tournaments on $14t$ candidates with block sizes $(7t,7t)$ and $(t,9t,3t,t)$. Then
\[
 \lin(A_t)-\lin(B_t)=6t^3,\qquad c(B_t)-c(A_t)=3t^3,\qquad
 \log\frac{p(B_t)}{p(A_t)}=3(\log 2)t^2+O(t\log(t+1)).
\]
In particular $B_t$ is more cyclic and, for all sufficiently large odd $t$, more probable, contradicting the conjecture.
\end{theorem}

\begin{proof}
The moments of the two partitions are $\sum_j b_j^2=98t^2,\ 92t^2$ and $\sum_j b_j^3=686t^3,\ 758t^3$ for $A_t,B_t$ respectively. Equation~\eqref{eq:blocklin} gives the linearity difference $(758-686)t^3/12=6t^3$, and~\eqref{eq:triangles} the cyclic-triple difference. Since $m=(\sum_j b_j^2-n)/2$, equation~\eqref{eq:blockprob} gives the probability ratio.
\end{proof}

\begin{remark}\label{rem:headline-scope}
The error term in Theorem~\ref{thm:headline} has an unspecified constant, so the theorem by itself certifies no particular $t$; Section~\ref{sec:finite} supplies explicit examples. For each fixed pair with $p(B)>p(A)$, the strict inequality also holds for all sufficiently large odd finite electorates, by convergence of both probabilities; the proof gives no electorate size that is uniform in $t$.
\end{remark}

The conjecture compares labelled tournaments. A natural variant compares the probabilities of isomorphism classes, $P[T]=\frac{n!}{|\mathrm{Aut}(T)|}p(T)$. The family also disproves this variant: $\log(n!/|\mathrm{Aut}(T)|)$ lies between $0$ and $\log n!=O(t\log t)$ for $n=14t$, which is absorbed in the error term of Theorem~\ref{thm:headline}.

\section{An explicit certified example and small cases}\label{sec:finite}

This section gives a fully explicit counterexample with a rigorous rational certificate, extends it to isomorphism classes, and reports numerical evidence on how small a counterexample can be.

\subsection{A certificate on 78 candidates}

In this subsection blocks are listed from \emph{weakest to strongest}, so that each later block beats every earlier block (the reverse of Section~\ref{sec:blocks}; reversing all arcs and the block order does not change probabilities, because $X$ and $-X$ have the same law).

\begin{proposition}\label{prop:78}
Let $A$ be an ordered regular block tournament on 78 candidates with two blocks of size 39, each carrying the triangle-decomposed tournament described below, and let $B$ be the ordered block tournament with circulant blocks of sizes $(7,7,49,5,5,5)$. Then $\lin(A)=145275$, $\lin(B)=145269$, $c(A)=4940$, $c(B)=4943$, and
\[
 p(B)/p(A)>10^5 .
\]
\end{proposition}

\begin{proof}
The linearities follow from~\eqref{eq:blocklin} and the cyclic-triple counts from~\eqref{eq:triangles}.

\emph{The blocks of $A$.} Label a block by $(x,i)\in\mathbb Z_{13}\times\mathbb Z_3$. Use the 13 triples $((x,0),(x,1),(x,2))$ and the 234 triples
\[
 ((x,i),(y,i),(7(x+y),i+1)),\qquad 0\le x<y\le12,\quad i\in\mathbb Z_3 ,
\]
each oriented cyclically in the displayed order. Each pair in a common layer occurs once. For a cross-layer pair $((u,i),(v,i+1))$, either $u=v$ (a vertical triple) or the other member of layer $i$ is uniquely $2v-u$ modulo 13, since $7$ is the inverse of $2$. Thus the 247 directed triangles partition the 741 pairs, and each vertex wins once in each of its 19 triangles, so the tournament is regular.

\emph{A sharper concavity bound.} For $t\le0$, let $Y=t-Z$ conditional on $Z\le t$, where $Z$ is standard normal. Its density $g$ on $y\ge0$ is proportional to $\exp(ty-y^2/2)$ and is decreasing. Hence $Y$ has the law of $UV$ with $U$ uniform on $[0,1]$ independent of $V$, where $V$ has density $-vg'(v)$ (integration by parts), and
\[
 \E[(Y-\E Y)^3]=\tfrac14\E[V(V-\E V)^2]\ge0 .
\]
In this exponential family the derivative of $\operatorname{Var}Y$ in $t$ is the third central moment, so $\operatorname{Var}Y\le 1-2/\pi$, its half-normal value at $t=0$. Since $1+\ell''(t)=\operatorname{Var}Y$ (proof of Theorem~\ref{thm:mode}), $\ell''(t)\le-2/\pi$ for $t\le0$. With concavity on the positive half-line this gives, for $z_-=\max(-z,0)$,
\[
 \ell(z)\le -\log2+r(0)z-z_-^2/\pi .
\]
Along a directed triangle the three strength differences sum to zero, so the sum of their squared positive parts is at most the square of the total negative magnitude, which is at most twice the sum of squared negative parts; hence the squared negative parts sum to at least one third of the squared norm of the three differences. The linear terms cancel on every directed triangle. Writing $m_A$ for the number of internal pairs of $A$, the internal likelihood is therefore at most $2^{-m_A}\exp(-\sum_{\rm int}(X_i-X_j)^2/(3\pi))$.

\emph{Upper bound for $A$.} Bound every cross-block term by the tangent of $\ell$ at $h=229/100$. The cross-block linear term is $r(h)a^2$ times the difference of the two block averages, which is centred normal with variance $2/a$ and independent of the within-block deviations ($a=39$). Within each block the sum of pairwise squared differences is $a$ times a chi-squared variable with $a-1$ degrees of freedom. Integrating these factors,
\begin{equation}\label{eq:upper78}
 p(A)\le 2^{-m_A}\left(1+\frac{2a}{3\pi}\right)^{-(a-1)}
             \Phi(h)^{a^2}\exp\{-a^2h r(h)+a^3r(h)^2\}.
\end{equation}

\emph{Lower bound for $B$.} Give the successive blocks of $B$ the constant strengths
\[
 u=(-275,\,-183,\,-4,\,164,\,219,\,298)/100 .
\]
The lower bound of Theorem~\ref{thm:mode} at this point gives
\begin{equation}\label{eq:lower78}
 p(B)\ge
 2^{-m_B}79^{-77/2}\exp\!\left(-\tfrac12\sum_i b_i u_i^2\right)
 \prod_{i<j}\Phi(u_j-u_i)^{b_i b_j}.
\end{equation}
The ratio of the right side of~\eqref{eq:lower78} to that of~\eqref{eq:upper78} exceeds $10^5$; Appendix~\ref{app:certificate} describes the rational arithmetic, and the program \texttt{certificate.py} performs it exactly. The certified rational lower bound is about $1.10\times10^5$; the bounds, not the true ratio, set this figure.
\end{proof}

\begin{corollary}\label{cor:classes}
For the tournaments of Proposition~\ref{prop:78}, $|\mathrm{Aut}(A)|=9$ and $|\mathrm{Aut}(B)|=300125$, so the isomorphism class of $B$ is more than $3.3$ times as probable as that of $A$.
\end{corollary}

\begin{proof}
In an ordered block tournament every vertex of a block of size $b$ has score $s+(b-1)/2$, with $s$ the number of vertices in weaker blocks. These scores are $19,58$ in $A$ and $3,10,38,65,70,75$ in $B$, distinct across blocks, so every automorphism maps each block to itself and $|\mathrm{Aut}|$ is the product of the blocks' automorphism counts. The circulant tournament on $\mathbb Z_b$ has exactly $b$ automorphisms: rotations are automorphisms, and conversely the out-neighbourhood $u+1,\dots,u+(b-1)/2$ of each vertex induces a transitive tournament in this order, so an automorphism $f$ satisfies $f(u+i)=f(u)+i$ and is a rotation. This gives $|\mathrm{Aut}(B)|=7\cdot7\cdot49\cdot5^3=300125$. An exhaustive search in \texttt{certificate.py} finds exactly three automorphisms of each 39-vertex block of $A$. Since $P[T]=\frac{n!}{|\mathrm{Aut}(T)|}p(T)$, the class ratio is $\frac{9}{300125}\,p(B)/p(A)$, and the certified bound exceeds $3.3$.
\end{proof}

The same program checks a second, more elementary witness on 319 candidates, $(159,159,1)$ against $(61,197,61)$, with ratio above $2^{39}$, using only the bound $p\le 2^{-m}$ for regular blocks and Theorem~\ref{thm:mode}.

\subsection{How small can a counterexample be?}\label{sec:small}

The certificate's bounds lose a factor of order $(n+1)^{(n-1)/2}$, so they cannot reach small $n$. To locate the threshold we estimated $p(T)$ numerically. Each probability is the $n$-dimensional integral~\eqref{eq:conditional}; we used importance sampling from a Student-$t$ distribution (8 degrees of freedom) centred at the maximiser of $H_T$ with the inverse negative Hessian as scale (program \texttt{small\_reversals.py}). These are Monte Carlo estimates with standard errors, not certificates.

\begin{itemize}[leftmargin=*]
\item \emph{All tournaments on $n\le7$.} For each of the 2, 4, 12, 56 and 456 isomorphism types on $n=3,\dots,7$ vertices, we estimated $p(T)$ with $2\times10^5$ samples (relative standard error below $0.15\%$). No reversal occurs: whenever $\lin(T)>\lin(U)$, the estimated ratio $p(T)/p(U)$ is at least $1.29$ for $n=7$ and larger for smaller $n$. This agrees with Harrison-Trainor's values for $n\le5$ \cite{harrison}.
\item \emph{Block tournaments on $8\le n\le 14$.} Among all ordered regular block tournaments with circulant blocks (21 to 377 block patterns for each $n$; $4\times10^4$ samples each), the largest estimated ratio of a more cyclic to a less cyclic one is below one for every $n\le12$ (it is about $0.91$ at $n=12$) and exceeds one at $n=13$ and $n=14$.
\item \emph{High-precision checks.} Table~\ref{tab:small} gives high-precision estimates for four pairs. On 14 candidates the block structure $(1,1,1,9,1,1)$ is about $2.7$ times as probable as $(7,7)$ although it has two more cyclic triples, and the $t=1$ member $(1,9,3,1)$ of Theorem~\ref{thm:headline} is already reversed, by a factor of about $1.27$. On 13 candidates two pairs are reversed by about $0.6\%$ and $1\%$; these margins are small but many times the spread between independent runs.
\end{itemize}

\begin{table}[htbp]
\centering
\begin{tabular}{llcccc}
\toprule
Less cyclic $U$ & More cyclic $T$ & $\lin(U),\lin(T)$ & $c(U),c(T)$ & $p(T)/p(U)$ & s.d.\ of log\\
\midrule
$(7,7)$ & $(1,1,1,9,1,1)$ & $763,\ 759$ & $28,\ 30$ & $2.706$ & $0.0021$\\
$(7,7)$ & $(1,9,3,1)$ & $763,\ 757$ & $28,\ 31$ & $1.268$ & $0.0022$\\
$(7,1,1,1,3)$ & $(3,7,3)$ & $620,\ 618$ & $15,\ 16$ & $1.0058$ & $0.0009$\\
$(7,1,1,1,1,1,1)$ & $(1,1,7,3,1)$ & $622,\ 620$ & $14,\ 15$ & $1.0095$ & $0.0005$\\
\bottomrule
\end{tabular}
\caption{Monte Carlo estimates (four independent runs of $10^6$ importance samples per tournament; the last column is the standard deviation of the log ratio across the four runs). Block sizes are listed from strongest to weakest. These numbers are numerical evidence, not certificates.}\label{tab:small}
\end{table}

Together with Proposition~\ref{prop:78}, this suggests that the smallest counterexample has between 8 and 13 candidates. For $8\le n\le12$ only block tournaments were searched. Certifying a 13-candidate reversal would need a 13-dimensional integral evaluated rigorously to better than $0.5\%$, which we have not attempted.

\section{The transitive tournament maximises probability}\label{sec:compression}

This section proves the half of the conjecture that survives: linearity does identify the most probable tournaments. The argument uses only the ordering of the latent strengths, so it applies to a wide class of paired-comparison models.

Let $X_1,\ldots,X_n$ be exchangeable real strengths with $\Pr(X_i\ne X_j)>0$ for $i\ne j$. Given the strengths, orient the pairs independently, $i$ beating $j$ with probability $K(X_i,X_j)$, where $0<K<1$, $K(x,y)+K(y,x)=1$, and $K(x,z)$ is strictly increasing in $x$ for every $z$. These conditions are a form of strong stochastic transitivity \cite{sst} and include the election model $K(x,y)=\Phi(x-y)$ of Section~\ref{sec:gauss}. Write $\Pr(T)$ for the resulting probability of a labelled tournament.

Given an arc $i\to j$ of $T$, the \emph{compression} $C_{ij}T$ replaces every path $j\to k\to i$ by $i\to k\to j$ and keeps all other arcs.

\begin{theorem}\label{thm:compression}
If $C_{ij}T\ne T$, then $\Pr(C_{ij}T)>\Pr(T)$. Consequently the labelled tournaments of maximum probability are exactly the transitive tournaments.
\end{theorem}

\begin{proof}
Condition on the strengths of all vertices other than $i,j$ and on the unordered pair $\{X_i,X_j\}$. By exchangeability the two assignments of that pair to $i$ and $j$ are equally likely. If the two values coincide, both tournaments have the same conditional likelihood. Otherwise write $h>l$ for the two values, $q=K(h,l)>1/2$ and $a_0=\frac12\log(q/(1-q))>0$.

For each vertex $k$ whose comparisons with $i$ and $j$ have different outcomes, let $z=X_k$ and put
\[
 U_k=K(h,z)[1-K(l,z)],\quad V_k=K(l,z)[1-K(h,z)],\quad
 a_k=\tfrac12\log(U_k/V_k)>0 .
\]
Write $s_k=1$ for the path $i\to k\to j$ and $s_k=-1$ for $j\to k\to i$. Arcs not meeting $\{i,j\}$, and the comparisons of $i$ and $j$ against a common opponent with identical outcomes, contribute a common factor to both assignments and both tournaments. Multiplying the remaining factors and averaging over the two assignments of $\{h,l\}$ gives
\[
 D\cosh\!\Big(a_0+\sum_k s_k a_k\Big)
\]
with a positive common $D$ containing $\sqrt{q(1-q)}\prod_k\sqrt{U_kV_k}$. Compression makes every $s_k$ positive. If some $s_k=-1$, strict positivity of $a_0$ and of the $a_k$ gives $|a_0+\sum_k s_ka_k|<a_0+\sum_k a_k$, so the conditional likelihood strictly increases whenever the two strengths differ, an event of positive probability.

Every nontransitive tournament has a directed triangle; compress one of its arcs. If $r,s$ count the paths $i\to k\to j$ and $j\to k\to i$ before compression, the move changes $d_i,d_j$ by $+s,-s$ and keeps the other scores, and as $d_i-d_j=1+r-s$ the linearity increases by $2s(r+1)>0$. As there are finitely many tournaments, repeated compression reaches a transitive tournament with strictly increasing probability, and exchangeability gives all transitive labellings the same probability.
\end{proof}

A related majorisation fact, that the score vector of a transitive tournament majorises every score vector, is used by Saile and Suksompong \cite[Lemma~6]{saile} for extremal probabilities of tournament solutions under independent edges; Theorem~\ref{thm:compression} instead concerns the probability of a single labelled tournament after averaging over exchangeable strengths. An exhaustive check over all 1024 labelled tournaments on five vertices in a two-strength model (\texttt{check\_compression.py}) confirms that all 5920 nontrivial compressions increase probability and that exactly the 120 transitive tournaments attain the maximum.

\part*{Part II. What the win totals of a round robin reveal}
\addcontentsline{toc}{part}{Part II}

The conjecture asked whether the win totals (the score vector) of a majority tournament determine its probability ordering. Part~II asks the same question for general paired-comparison models, where it becomes a question about goodness of fit: in a single round robin, such as a sports league or a set of pairwise judgements, how much information about the comparison model remains once each player's win total is known? Part~I does not depend on Part~II.

\section{Exponential separation at fixed win totals}\label{sec:general-kernels}

This section shows that, outside the Bradley--Terry family, typical outcomes have same-score neighbours whose probabilities differ by factors $e^{\pm cn^2}$.

\paragraph{Setting.} Let $(\mathcal X,\mu)$ be a standard probability space and $W$ a measurable \emph{tournament kernel}: $W(x,y)+W(y,x)=1$. Draw latent types $X_1,\dots,X_n$ independently from $\mu$ and, given the types, orient the pairs independently with $\Pr(i\to j)=W(X_i,X_j)$. Write $Q_x$ for the conditional law of the tournament given types $x$, $P_n$ for its mixture over types, $d(T)=(d_1,\dots,d_n)$ for the labelled score vector, $\mathcal F_d=\{T:d(T)=d\}$ for a score class, and $U_d$ for the uniform law on $\mathcal F_d$. A law is \emph{uniform given the scores} if its conditional law on each score class of positive mass is $U_d$. We call $W$ \emph{interior} if $0<W<1$ $\mu^2$-almost everywhere. For three types write $q_\pm$ for the probabilities of the two directed triangles,
\[
 q_+=W(x,y)W(y,z)W(z,x),\qquad q_-=W(y,x)W(z,y)W(x,z),
\]
and call $W$ \emph{triangle balanced} if $q_+=q_-$ almost everywhere. The Bradley--Terry (logistic) kernels $W(x,y)=(1+e^{a(y)-a(x)})^{-1}$ are triangle balanced. For example, the election model of Part~I has $W(x,y)=\Phi(x-y)$, which is not.

\paragraph{What was known.} The qualitative statement that score-uniformity forces the Bradley--Terry form follows from Lauritzen's theorem on exchangeable Rasch arrays: a separately exchangeable binary array whose submatrix probabilities depend only on row and column sums has a kernel of ``Rasch type'' \cite[Section~2.4, eq.~(2)]{Lauritzen2008}, a result Lauritzen attributes to his 2003 work. Splitting the vertices of a score-uniform tournament law into two groups and summing over within-group outcomes produces such an array, and reciprocity turns the Rasch form into a difference of vertex potentials. For independent edges with fixed probabilities on a finite tournament, Iyer proves the equivalence of uniformity given the scores and the Bradley--Terry form by triangle reversals \cite[Theorem~4.2.9]{Iyer2016}. Georgakopoulos and Tyros characterise Bradley--Terry probabilities by balance on triangles and test this property with query access to the probabilities \cite{GT2016}, and Th\"ornblad studies expected score functions of tournament limits \cite{Thornblad2016}. The new content below is quantitative: it concerns the realised score class of a typical outcome after averaging over latent types.

\begin{theorem}\label{thm:general-kernels}
If $W(x,y)=(1+\exp(a(y)-a(x)))^{-1}$ $\mu^2$-almost everywhere for a finite measurable $a$, then $P_n$ is uniform given the scores for every $n$. For every other interior kernel $W$ there is $c>0$ such that:
\begin{enumerate}
\item with $P_n$-probability tending to one, $T$ has tournaments $T_-,T_+$ with exactly its labelled scores and
\[
 P_n(T_-)\le e^{-cn^2}P_n(T),\qquad P_n(T_+)\ge e^{cn^2}P_n(T);
\]
\item there are events $A_n$ with $P_n(A_n)\to1$ and $\sup_d U_d(A_n)\le e^{-cn^2}$ for all large $n$; hence for every sequence $R_n$ of laws uniform given the scores,
\[
 \|P_n-R_n\|_{\rm TV}\to1\quad\text{and}\quad\liminf_{n\to\infty} n^{-2}D_{\rm KL}(P_n\Vert R_n)>0.
\]
\end{enumerate}
\end{theorem}

\begin{proof}
In the Bradley--Terry case put $w_i=e^{a(x_i)}$. Conditional on the types, $Q_x(T)=\prod_i w_i^{d_i(T)}/\prod_{i<j}(w_i+w_j)$, which depends on $T$ only through $d(T)$; averaging over $x$ preserves this.

For the converse let $H(x,y)=\log[W(x,y)/W(y,x)]$ and $D=\log(q_+/q_-)=H(x,y)+H(y,z)+H(z,x)$. If $D=0$ almost everywhere, Fubini's theorem gives a fixed $z$ for which this holds for almost every $(x,y)$ with $H(x,z)$ finite, and $a(x)=H(x,z)$ gives the Bradley--Terry form after redefinition on a null set. Otherwise $D>0$ on a set of positive measure, since swapping two arguments reverses its sign.

We use twice the following observation: a bounded measurable $f$ on $\mathcal X^3$ that is positive on a set of positive measure has strictly positive integral on some product $B_1\times B_2\times B_3$. Indeed $\{f>0\}$ can be approximated in measure by a finite union of measurable rectangles, and partitioning that union into disjoint rectangles keeps one positive integral.

\emph{Triangle packing.} Put $m=\lfloor n/3\rfloor$ and, before observing anything, fix three disjoint groups of $m$ labels, $A_i,B_i,C_i$ with indices modulo $m$, and the $m^2$ triples $(A_i,B_j,C_{i+j})$. Each pair from different groups lies in exactly one triple, so the triples are edge-disjoint, and each vertex lies in $m$ of them. The forward orientation is $A_i\to B_j\to C_{i+j}\to A_i$. A \emph{label selection} $b$ is a choice of a subset of each group; a triple is selected if all three of its vertices are chosen. For a fixed rectangle $B_1\times B_2\times B_3$, choosing vertices by their latent types defines a random selection among the $2^{3m}$ label selections. For any bounded summand depending on a triple's types and outcomes, its average over the packing converges in probability to its expectation, because vertex-disjoint triples are independent and each triple meets at most $3m$ others, so the variance of the average is $O(1/m)$.

\emph{The event $A_n$ and the less probable neighbour.} Apply the rectangle observation to $f=q_+-q_-$, with positive integral $\gamma$ on some rectangle. Let $S_b$ be the number of selected forward cycles minus selected reverse cycles. For the latent selection $S_b/m^2\to\gamma$, so the observable event $A_n=\{\max_b S_b\ge\gamma m^2/2\}$ has probability tending to one. Under $U_d$, fix $b$ and condition on the arcs outside the packing, the transitive packed triples and the positions of the cyclic packed triples. The remaining cyclic orientations are independent fair coins, since reversing any subset of them preserves the scores and the conditioning data. Hoeffding's inequality and a union bound give
\begin{equation}\label{eq:fair-cycle-tail}
 \sup_d U_d(A_n)\le 2^{3m}\exp(-\gamma^2m^2/8).
\end{equation}
Hence for some $a>0$ and all large $n$ every law $R_n$ uniform given the scores has $R_n(A_n)\le e^{-a n^2}$. Total variation is at least $P_n(A_n)-R_n(A_n)$, and the log-sum inequality on $A_n$ and its complement gives $D_{\rm KL}(P_n\Vert R_n)\ge P_n(A_n)\log(1/R_n(A_n))-\log 2$. Now take $R_n(T)=P_n(\mathcal F_{d(T)})/|\mathcal F_{d(T)}|$. For $0<b'<a$, the $P_n$-mass of $A_n\cap\{P_n(T)<e^{b'n^2} R_n(T)\}$ is at most $e^{b'n^2}R_n(A_n)\to0$. So with probability tending to one the sampled tournament exceeds its score-class average by a factor $e^{b'n^2}$, and some member $T_-$ of the class has probability at most the average.

\emph{The more probable neighbour.} For a label selection $b$ and observed $T$, flip an independent fair coin on each selected reverse cycle and reverse it on heads; call the output $T'_b$, which has scores $d(T)$. Since the selected triangles are edge-disjoint,
\[
 \E_{\rm flip}\frac{Q_x(T'_b)}{Q_x(T)}
 =\prod_{t\ \text{selected reverse cycle}}\frac{1+e^{D_t}}2 .
\]
Each factor has logarithm at least $-\log2$. Thus $f=q_-\min\{1,\log((1+e^D)/2)\}$ is bounded and positive where $D>0$; let $\eta>0$ be its integral on some rectangle. For the corresponding latent selection, the packing average of selected reverse-cycle indicators times their clipped log factors converges to $\eta$, so with probability tending to one
\[
 \max_b\E_{\rm flip}\frac{Q_X(T'_b)}{Q_X(T)}\ge\exp(\eta m^2/2).
\]
Call this event $E$; with probability tending to one, $\Pr(E\mid T)>1/2$. For each fixed selection the flipping rule depends only on $T$ and fresh coins, so Bayes' identity gives
\[
 \sum_b\E_{\rm flip}\frac{P_n(T'_b)}{P_n(T)}
 =\E\Big[\sum_b\E_{\rm flip}\frac{Q_X(T'_b)}{Q_X(T)}\,\Big|\,T\Big]
 \ge\Pr(E\mid T)e^{\eta m^2/2}.
\]
With $2^{3m}$ selections, some selection and some realisation of its flips give a ratio at least $2^{-3m-1}e^{\eta m^2/2}$. This proves the claim for $T_+$ after reducing $c$.
\end{proof}

The proof allows interior probabilities arbitrarily close to zero or one and needs no moment condition on the log odds. It also covers kernels in which every type has the same expected win fraction, so that the expected scores carry no information at all.

\begin{corollary}[Kernels with deterministic comparisons]\label{cor:boundary}
A $[0,1]$-valued tournament kernel is uniform given the scores for all $n$ if and only if it is triangle balanced. Every other such kernel satisfies part~2 of Theorem~\ref{thm:general-kernels} and the existence of $T_-$ in part~1.
\end{corollary}
\begin{proof}
Under balance, almost every type vector has $q_+=q_-$ on every triple, so reversing a directed triangle preserves the conditional probability. Any two tournaments with the same labelled scores are connected by triangle reversals: the arcs where they disagree, oriented as in the first, have equal in- and outdegree at each vertex and decompose into directed cycles; a cycle of length above three is split at a chord by reversing the shorter cycle through the chord and then the other, which uses the reversed chord. Induction on length implements both by triangle reversals and restores the chord. Hence all probabilities on a score class agree, including zeros. In the unbalanced case the argument with $q_+-q_-$ above applies verbatim.
\end{proof}

The more probable neighbour can fail at the boundary. Give each vertex a uniform group in $\{0,1,2\}$ and an independent uniform rank; between groups use the deterministic cycle $0\to1\to2\to0$, and within groups order by rank. Three distinct groups violate triangle balance, but at most $3^n n!$ tournaments can occur, so for each $b>0$ the outputs of probability at most $e^{-bn^2}$ have total mass at most $3^n n!e^{-bn^2}\to0$; with probability tending to one even the most probable output has ratio below $e^{bn^2}$ to the sampled one.

\paragraph{Scalar links.} Suppose $W(x,y)=G(x-y)$ on the real line, with $G$ continuous, strictly increasing and $G(-t)=1-G(t)$, and $\mu$ of full support.

\begin{corollary}\label{thm:score-link}
Such a model is uniform given the scores for every $n$ exactly when $G(t)=(1+e^{-bt})^{-1}$ for some $b>0$. Every other such link exhibits both exponential gaps of Theorem~\ref{thm:general-kernels}; in particular the election model $G=\Phi$ does.
\end{corollary}
\begin{proof}
Write $L(t)=\log[G(t)/G(-t)]$. By continuity and full support, triangle balance at every triple of strengths is $L(u)+L(v)=L(u+v)$; continuous additive functions are linear, and strict increase gives positive slope. The converse is the Bradley--Terry case. For $G=\Phi$, $1-\Phi(t)\le e^{-t^2/2}$ gives $L(t)\ge t^2/2-\log 2$, which is not linear.
\end{proof}

\section{Conditional tests from one round robin}\label{sec:tests}

This section turns Theorem~\ref{thm:general-kernels} into tests that use only the observed tournament, are exactly calibrated given the win totals, and detect fixed alternatives. The null hypothesis is ``uniform given the scores'', which by Section~\ref{sec:general-kernels} is the Bradley--Terry model with arbitrary, possibly random, strengths; conditioning on the scores removes the nuisance strengths. All tests use the labelled triangle packing of Section~\ref{sec:general-kernels}, chosen before the data are seen, and the sign $Y_t\in\{+1,-1,0\}$ of each packed triangle $t=(A_i,B_j,C_{i+j})$: $+1$ for the forward cycle, $-1$ for the reverse cycle, $0$ if transitive. Reversing any set of cyclic packed triangles preserves all scores, so under every law uniform given the scores the cyclic signs are independent fair signs given the scores, the other arcs and the cycle positions. This exact symmetry is the source of calibration in every test below; \texttt{check\_cycles.py} confirms it exhaustively for all 33,864 labelled tournaments on three to six vertices. Classical lack-of-fit tests also count cyclic triads, and Singh and Davidov \cite{SD2026} develop more powerful lack-of-fit tests based on cyclic structure for cardinal (continuous) paired-comparison data on dense and sparse comparison graphs, and discuss extensions to binary data. The tests below differ in setting and guarantee: they use binary outcomes from a single round robin, are exactly calibrated on every score class, and their reported significance decays at a quadratic-exponential rate under each fixed alternative.

\subsection{A spectral test for an arbitrary kernel}\label{sec:kernel-test}

Before inspecting results, put each unordered pair independently into a training set with probability $1/2$. Let $D_{ij}=2$ if the pair is in the training set and $i$ wins, $D_{ij}=-2$ if it is in the training set and $j$ wins, and $D_{ij}=0$ otherwise. Compute a best rank-$r$ approximation $D_r$ with $r=\lfloor\sqrt n\rfloor$ (ties broken by a fixed rule), project it onto the skew-symmetric matrices, clip entries to $[-1,1]$ to obtain $\widehat M$, and put $\widehat p_{ij}=(1+\widehat M_{ij})/2$. This is a rank-truncation form of spectral matrix estimation in the setting of Chatterjee \cite{usvt}. Retain the packed triangles whose three pairs were all held out, and for retained $t=(a,b,c)$ put
\[
 f(p,q,s)=pqs-(1-p)(1-q)(1-s),\qquad
 w_t=f(\widehat p_{ab},\widehat p_{bc},\widehat p_{ca}),
\]
$S=\sum_t w_tY_t$ and $V=\sum_t w_t^2|Y_t|$. Report $p_{\rm sp}=\exp[-(\max(S,0))^2/(2V)]$ if $V>0$ and $1$ otherwise.

\begin{proposition}\label{prop:kernel-test}
Conditional on every fixed training mask and every feasible score vector, under the uniform law on that score class, $\Pr(p_{\rm sp}\le\alpha)\le\alpha$ for $0<\alpha<1$. For a fixed $[0,1]$-valued kernel $W$ with $A=\E[(q_+-q_-)^2]>0$, where $q_\pm$ are evaluated at three independent types, put $B=\E[(q_+-q_-)^2(q_++q_-)]$. Then $-n^{-2}\log p_{\rm sp}\to A^2/(144B)\ge A/144$ in probability, including the randomness of the mask. The statistic takes $O(n^3)$ operations with a dense singular value decomposition.
\end{proposition}

\begin{proof}
\emph{Size.} Condition further on every arc outside the retained packing, the transitive retained triangles and the positions of the cyclic ones. The training entries, and hence the weights, are fixed, and the cyclic signs are independent fair signs. Their weighted sum has conditional upper tail at most $\exp(-s^2/(2V))$ for $s>0$.

\emph{Estimation error.} Conditional on the types, the mean $M$ of $D$ has entries $M_{ij}=2W(X_i,X_j)-1$ off the diagonal, and $E=D-M$ is skew-symmetric with independent bounded centred entries above the diagonal. For unit vectors $u,v$ the coefficients of these entries in $u^{\mathsf T}Ev$ have squared sum $\sum_{i<j}(u_iv_j-u_jv_i)^2=1-(u^{\mathsf T}v)^2\le1$, so Hoeffding's inequality and a union bound over two $1/4$-nets of at most $9^n$ vectors give $\|E\|_{\rm op}=O_P(\sqrt n)$ uniformly in the types. Comparing $\|D-D_r\|_{\rm F}$ with $\|D-M_r\|_{\rm F}$ for a best rank-$r$ approximation $M_r$ of $M$, and using the nuclear-norm bound for a matrix of rank at most $2r$,
\[
 \|D_r-M\|_{\rm F}^2-\|M_r-M\|_{\rm F}^2
 \le 2\|E\|_{\rm op}\sqrt{2r}\,\|D_r-M_r\|_{\rm F},
\]
whence $\|D_r-M\|_{\rm F}\le\|M_r-M\|_{\rm F}+2\sqrt{2r}\|E\|_{\rm op}$. For each $\eta>0$, approximate $2W-1$ in $L^2(\mu^2)$ within $\eta$ by a bounded step kernel on a finite measurable partition; its sampled matrix has bounded rank and expected squared Frobenius error at most $(\eta^2+O(1/n))n^2$. Markov's inequality and $\eta\downarrow0$ give $\|M_r-M\|_{\rm F}/n\to0$, and projection and clipping do not increase the error, so $\|\widehat M-M\|_{\rm F}/n\to0$ in probability.

\emph{Limits.} For a packed triangle put $F_t=f(W(X_a,X_b),W(X_b,X_c),W(X_c,X_a))=q_+-q_-$. Each partial derivative of $f$ has magnitude at most one on the unit cube and no two packed triangles share a pair, so Cauchy--Schwarz gives $m^{-2}\sum_{t}|w_t-F_t|=o_P(1)$ with $m=\lfloor n/3\rfloor$. Given the types, $\E[Y_t\mid X]=F_t$, and each triangle is retained with probability $1/8$ independently of its types and outcomes, and its weight is independent of its own outcome. Vertex-disjoint terms are independent, so the variance of the sum is $O(m^3)$, and $S/m^2\to\frac18\E[F_t^2]=A/8$ in probability. Similarly $|w_t^2-F_t^2|\le2|w_t-F_t|$ and $\E[|Y_t|\mid X]=q_++q_-$ give $V/m^2\to B/8$. Here $B>0$ because $q_++q_-\ge|q_+-q_-|$, and $B\le A$ because the two directed triangles are disjoint events, so $q_++q_-\le1$. Substituting the limits and $m/n\to1/3$ gives $-\log p_{\rm sp}/n^2\to (A/8)^2/(2\cdot 9\cdot B/8)=A^2/(144B)$.
\end{proof}

The test also bounds the information lost by keeping only the scores. Writing $H(T\mid d(T))$ for conditional entropy in nats, whenever $A>0$,
\begin{equation}\label{eq:entropy-deficit}
 \liminf_{n\to\infty}
 \frac{\E_{P_n}\log |\mathcal F_{d(T)}|-H_{P_n}(T\mid d(T))}{n^2}
 \ \ge\ \frac{A^2}{144B}\ \ge\ \frac{A}{144}.
\end{equation}
Indeed, let $R_n$ be uniform given the scores with the same score law as $P_n$; the numerator is $D_{\rm KL}(P_n\Vert R_n)$, which is unchanged by adjoining the same independent training mask to both laws. For $0<s<A^2/(144B)$ the event $p_{\rm sp}\le e^{-sn^2}$ has probability tending to one under $P_n$ and at most $e^{-sn^2}$ under $R_n$; apply the log-sum inequality from the proof of Theorem~\ref{thm:general-kernels} and let $s\uparrow A^2/(144B)$.

\subsection{A score-localised scan for scalar links}\label{sec:adaptive}

For the scalar links of Corollary~\ref{thm:score-link} a simpler test needs no matrix estimation. At resolution $\rho=0,\ldots,\lfloor\log_2 n\rfloor$ put $q=2^\rho$ and $h_{\rho,b}(z)=(1-|qz-b|)_+$ for $b=0,\ldots,q$. For $\mathbf b=(b_1,b_2,b_3)$, let $w_t$ be the product of the three hat functions evaluated at the observed win fractions $d_v/(n-1)$ of the vertices of $t$, in group order, and set
\[
 S_{\rho,\mathbf b}=\sum_t w_t Y_t,\quad
 V_{\rho,\mathbf b}=\sum_t w_t^2 |Y_t|,\quad
 \alpha_{\rho,\mathbf b}=2^{-\rho-1}(q+1)^{-3},
\]
\begin{equation}
 p_{\rm scan}=\min\left\{1,\
 \min_{\rho,\mathbf b:\,V_{\rho,\mathbf b}>0}\frac{2}{\alpha_{\rho,\mathbf b}}
       \exp\left(-\frac{S_{\rho,\mathbf b}^2}{2V_{\rm \rho,\mathbf b}}\right)\right\}.
 \label{eq:scan}
\end{equation}
At most two hat functions per coordinate are positive at each resolution, so the scan takes $O(n^2\log n)$ expected hash-table operations.

\begin{corollary}\label{cor:scan}
For every feasible score vector $d$, $U_d(p_{\rm scan}\le\alpha)\le\alpha$ for $0<\alpha<1$. Under any fixed non-logistic link of Corollary~\ref{thm:score-link} there is $c>0$ with $P_n(p_{\rm scan}\le e^{-cn^2})\to1$.
\end{corollary}

\begin{proof}
\emph{Size.} The weights depend only on the scores, which reversals preserve, so given the conditioning data of the size argument above the cyclic signs are fair and the weights fixed. The bound $\cosh z\le e^{z^2/2}$ gives the two-sided tail $2\exp(-s^2/(2V_{\rho,\mathbf b}))$ in each cell, and the weights $\alpha_{\rho,\mathbf b}$ sum to at most one, so a union bound gives validity. (Enumerating the fair-sign assignments gives an exact randomisation $p$-value, which is at most~\eqref{eq:scan}.)

\emph{Power.} Let $L(t)=\log[G(t)/G(-t)]$ and choose $u,v>0$ with $L(u)+L(v)\ne L(u+v)$. At strength levels $u+v,v,0$ one directed cycle, call it $+$, is more probable than the other, and by continuity there are disjoint open strength intervals around the three levels on whose product the cycle odds exceed $e^h$ for some $h>0$. The mean win probability $s(x)=\int G(x-y)\,\dd\mu(y)$ is continuous and strictly increasing, and given $X_i$ the win fraction $r_i=d_i/(n-1)$ is an average of independent Bernoulli variables of mean $s(X_i)$, so Hoeffding's inequality gives
\begin{equation}\label{eq:score-bands}
 P\Big(\max_i|r_i-s(X_i)|>\delta\Big)\le 2n\exp\{-2(n-1)\delta^2\}.
\end{equation}
Fix a resolution fine enough that one product of hats has support inside the product of the images under $s$ of the three strength intervals. Evaluated at the latent scores $s(X_i)$, the expected weighted signed count per packed triangle in this cell is a fixed nonzero number, because the two cycle probabilities differ in the same direction throughout the box; vertex-disjoint triples are independent, so the variance of the sum is $O(m^3)$. Replacing latent scores by observed fractions changes the average by $o_P(1)$, by~\eqref{eq:score-bands} and the Lipschitz constant of the hats. Hence $|S_{\rho,\mathbf b}|/m^2$ converges to a positive number, while $V_{\rho,\mathbf b}\le m^2$ and $\alpha_{\rho,\mathbf b}$ is fixed.
\end{proof}

\subsection{A test targeted to the election model and its finite-sample behaviour}\label{sec:targeted}

When the alternative is the election law itself (probit link, Gaussian strengths), the expected direction of each cycle can be predicted from the scores. Under that law $s(x)=\Phi(x/\sqrt2)$; estimate strengths by $\widehat x_i=\sqrt2\,\Phi^{-1}((d_i+1/2)/n)$, put $L(z)=\log(\Phi(z)/\Phi(-z))$, and for a packed triangle $t=(a,b,c)$ let $D_t=L(\widehat x_a-\widehat x_b)+L(\widehat x_b-\widehat x_c)+L(\widehat x_c-\widehat x_a)$ and $u_t=\max(-4,\min(4,D_t))$. The statistic
\begin{equation}\label{eq:targeted}
 E_{\rm G}=\prod_{t:\,Y_t\ne0}\frac{2}{1+\exp(-Y_tu_t)}
\end{equation}
has conditional mean one under every law uniform given the scores, because the $u_t$ depend only on scores and each factor has mean one under a fair sign. So $\min(1,1/E_{\rm G})$ is a valid $p$-value (an e-value test). Under the election law $m^{-2}\log E_{\rm G}$ converges in probability to a positive constant: with true strengths, the conditional expected log factor of a cyclic triple is strictly concave in the proposed log odds, zero at zero and maximal at the true value, so clipping towards zero keeps a positive gain wherever the true cycle log odds are nonzero, which happens with positive probability; restricting to strengths in $[-K,K]$, \eqref{eq:score-bands} and continuity of $\Phi^{-1}$ make the estimated factors uniformly consistent, and $K\to\infty$ removes the restriction.

\begin{table}[htbp]
\centering
\begin{tabular}{rrcc}
\toprule
Candidates & Strength s.d. & Logistic null & Probit alternative\\
\midrule
90  & 1.0 & 0/256 & 0/256\\
180 & 1.0 & 1/256 & 11/256\\
360 & 1.0 & 1/256 & 133/256\\
600 & 1.0 & 0/256 & 245/256\\
360 & 0.5 & 3/256 & 29/256\\
360 & 1.5 & 2/256 & 128/256\\
\bottomrule
\end{tabular}
\caption{Rejections at $E_{\rm G}\ge20$ in simulated round robins with one outcome per pair and independent Gaussian strengths (256 replicates per cell; \texttt{power\_study.py}). The logistic link was $(1+e^{-1.7z})^{-1}$ and the probit link $\Phi(z)$. Only the unit-variance probit design is the limiting impartial-culture election.}\label{tab:power}
\end{table}

Table~\ref{tab:power} shows what these asymptotic guarantees mean at finite size. The logistic and probit links are notoriously hard to distinguish, and one round robin needs several hundred players: at unit strength variance there were no rejections at $n=90$, $133/256$ ($52.0\%$, Wilson 95\% interval $45.8$--$58.0\%$) at $n=360$, and $245/256$ at $n=600$. Halving the strength spread reduces the power at $n=360$ to $11.3\%$. For leagues of a few dozen teams these tests are calibrated but have little power against links this close to logistic; they are better suited to large pairwise-comparison data sets or to alternatives with genuine cyclic structure.

\section{One observation per pair is not enough for minimax testing}\label{sec:onerep}

This section answers a remark of Makur and Singh about replication. They study minimax testing of the Bradley--Terry--Luce (BTL) model when each \emph{ordered} pair $(i,j)$ supplies $k$ independent Bernoulli observations of mean $p_{ij}$, with no reciprocity imposed on the alternative \cite{MS2025}. Their upper bound (Theorem~2 there) needs $k\ge\max\{2,c\,d_{\max}h\log n/\xi^4\}$, which is of order $n\log n$ on the complete graph, and after equation~(26) they note that it ``remains unclear'' whether $k\ge2$ is necessary, recalling that in a related two-sample problem Rastogi et al.\ showed it to be \cite[Proposition~4]{RBSS2022}. We show that $k=1$ is impossible at constant separation, adapting the Bernoulli masking construction of Rastogi et al.\ to the one-sample composite null, and give a separate statement for reciprocal round-robin data. Whether two (or any constant number of) observations per pair suffice remains open.

\begin{proposition}\label{thm:one-rep}
(a) On the complete graph with one observation per ordered pair, the Makur--Singh minimax sum of type I and type II errors at separation $1/64$, in their separation distance and with fixed admissible expansion and principal-ratio constants, is at least $1-\exp[-\binom n2/8]$.

(b) With one observation per unordered pair (a single round robin), the minimax sum of errors for testing BTL against moderately stochastically transitive reciprocal alternatives with $\inf_{B\in\mathrm{BTL}}\|P-B\|_{\rm F}/n\ge 1/1000$ is at least $1-\exp(-\lfloor n/3\rfloor^2/128)$ for $n\ge6$.
\end{proposition}

\begin{proof}
(a) Use Makur and Singh's convention that $p_{ij}$ is the probability that the second team wins. Let $P_{ii}=1/2$ and draw each $p_{ij}$, $i\ne j$, independently and uniformly from $\{1/4,3/4\}$. With one observation per ordered pair the mixture of the data laws is exactly the equal-strength BTL law $P_0$, a single null point: each observation is an independent fair coin. In their notation let $\pi$ be the stationary law of the canonical Markov chain, $\Pi=\operatorname{diag}(\pi)$ and $H=\Pi P+P\Pi-\mathbf1\pi^\top$; their separation distance \cite[eq.~(17)]{MS2025} is $\|H\|_{\rm F}/(n\|\pi\|_\infty)$. Their Proposition~5 with dynamic-range parameter $1/3$ gives $\pi_{\max}/\pi_{\min}\le9$ and, on the complete graph, a uniform positive lower bound on expansion, so every matrix in the construction satisfies their assumptions. Let $N=\binom n2$ and $J=\#\{i<j:p_{ij}=p_{ji}\}\sim\operatorname{Bin}(N,1/2)$, so $\Pr(J<N/4)\le e^{-N/8}$ by Hoeffding's inequality. For each counted pair $|H_{ij}+H_{ji}|=(\pi_i+\pi_j)|p_{ij}+p_{ji}-1|\ge \pi_{\min}$, so on $\{J\ge N/4\}$, using Cauchy--Schwarz on each counted pair,
\[
 \frac{\|H\|_{\rm F}}{n\|\pi\|_\infty}
 \ge \frac{\pi_{\min}}{n\pi_{\max}}\sqrt{J/2}
 \ge \frac{\sqrt{1-1/n}}{36}>\frac1{64}\qquad(n\ge2).
\]
Restricting the mixing law to an event of probability $1-\eta$ changes the data mixture by at most $\eta$ in total variation, so every test's type I error at $P_0$ plus its average type II error under the restricted mixture is at least $1-\eta$, with $\eta\le e^{-N/8}$.

(b) Now let $p_{ij}$ be the probability that $i$ wins, with $p_{ji}=1-p_{ij}$. Split the vertices into three tiers of sizes differing by at most one, with baseline BTL strengths $4,2,1$; keep $p_{ij}=1/2$ within tiers, and between tiers perturb each higher-tier win probability by an independent fair sign times $1/48$. One observation per unordered pair has, after mixing, exactly the baseline BTL law. The baseline probabilities for high--middle, middle--low and high--low comparisons are $2/3,2/3,4/5$. In tier order every realised matrix satisfies $p_{ik}\ge\min(p_{ij},p_{jk})$ whenever $i$ precedes $j$ and $j$ precedes $k$, which is moderate stochastic transitivity \cite[eq.~(6)]{RBSS2022}: for three distinct tiers the smallest high--low probability $4/5-1/48$ exceeds both adjacent-tier probabilities $2/3+1/48$, and if two consecutive vertices share a tier the minimum is at most $1/2\le p_{ik}$.

Take $m=\lfloor n/3\rfloor$ vertices from each tier and the $m^2$ edge-disjoint triangles of Section~\ref{sec:general-kernels}. Independently across triangles, the signs on the high--middle, middle--low and high--low pairs are $(+,+,-)$ with probability $1/8$, and then the probability triple is $q=(11/16,11/16,187/240)$. Every BTL triple $(x,y,z)$ lies in $\mathcal B=\{(x,y,z)\in[0,1]^3:xy(1-z)=(1-x)(1-y)z\}$. The polynomial $F(x,y,z)=xy(1-z)-(1-x)(1-y)z$ equals $869/30720$ at $q$, and its partial derivatives have magnitude at most one on the unit cube, so $\operatorname{dist}(q,\mathcal B)\ge 869/(30720\sqrt3)>16/1000$. By Hoeffding's inequality, except on an event of probability at most $\exp(-m^2/128)$, at least $m^2/16$ triangles have this sign pattern, and then $\|P-B\|_{\rm F}^2\ge (16/1000)^2m^2/16$ for every BTL matrix $B$, so $\|P-B\|_{\rm F}/n\ge (4/1000)m/n\ge 1/1000$ since $m/n\ge1/4$ for $n\ge6$. The restriction argument of part~(a) completes the proof.
\end{proof}

Under the stronger assumption of strong stochastic transitivity a positive result does exist with one observation per pair: Shah et al.'s least-squares estimator \cite[Theorem~1 and (5b)]{sst} satisfies $n^{-1}\|\widehat P-P\|_{\rm F}\le C\log n/\sqrt n$ with probability at least $1-c_1 e^{-c_2 n}$, so rejecting when the normalised distance from $\widehat P$ to the BTL class exceeds $C\log n/\sqrt n$ has vanishing errors against alternatives at distance $2C\log n/\sqrt n$ (that estimator is not known to be computable in polynomial time). Proposition~\ref{thm:one-rep}(b) shows that moderate stochastic transitivity is not enough. In practice this matters for any pairwise evaluation in which each pair is compared once, for example sports round robins or one-shot pairwise preference judgements: without replication, or structural assumptions such as those of Part~II, no test can reliably tell a BTL world from one far from it.

\section{Discussion}\label{sec:discussion}

This section summarises the results, their limits and their uses, and lists open problems.

\paragraph{Summary.} Under impartial culture, the probability of a large majority tournament is driven by how many pairs of candidates must end up nearly tied in strength, not by how many cyclic triples it contains. Ordered block structures make the two quantities disagree, so a more cyclic tournament can be exponentially more probable. The conjecture's extremal claim survives: the transitive tournaments are the most probable ones, for every $n$ and for a broad class of models. More generally, outside the Bradley--Terry family the win totals of a round robin leave a large, detectable amount of information about the comparison law.

\paragraph{Scope.} The results describe the shape of the impartial-culture distribution in the limit of many voters. They say nothing about real electorates, which are not impartial cultures, and the tournaments involved are individually very improbable (on 78 candidates, $p(A)$ is below $e^{-1000}$). They matter for exact and simulation work with the impartial-culture law, for example when tournaments are weighted or sampled by probability, and they explain why ``fewer cycles'' is not a reliable proxy for ``more likely''.

\paragraph{Uses of Part II.} The conditional tests apply to any single round robin: a sports league season, a tournament of game-playing programs, or a set of pairwise preference judgements in which each pair is judged once. Their calibration is exact for every score vector, but their power at realistic sizes is modest against near-logistic alternatives (Table~\ref{tab:power}). Proposition~\ref{thm:one-rep} gives practical advice for pairwise evaluation: compare each pair more than once, or accept a structural assumption, before testing the Bradley--Terry model.

\paragraph{Open problems.}
\begin{enumerate}[leftmargin=*]
\item Determine the smallest number of candidates admitting a reversal. The evidence of Section~\ref{sec:small} places it in $8\le n\le13$; a certificate at $n=13$ or $14$ would need validated integration in 13 or 14 dimensions.
\item Find a correct replacement for the conjecture. Theorem~\ref{thm:mode} shows that $M_T$ predicts $\log p(T)$ to within $O(n\log n)$; is there a combinatorial statistic with the same accuracy?
\item Do two, or any constant number of, observations per ordered pair suffice for the Makur--Singh guarantee?
\item Give finite-sample power bounds for the tests of Section~\ref{sec:tests} in terms of the triangle imbalance $A$ and the matrix-estimation error.
\end{enumerate}

\paragraph{AI-use statement.} This work was produced with the assistance of AI systems, which generated the mathematical arguments, computations and text. The proofs were checked by internal AI-assisted review and by the companion programs; no external peer review has taken place.

\appendix

\section{Rational arithmetic of the 78-candidate certificate}\label{app:certificate}

This appendix records how \texttt{certificate.py} compares the right sides of~\eqref{eq:lower78} and~\eqref{eq:upper78} without floating point. The ratio is squared so that every exponent on $\Phi$ and on the determinant factor is an integer. The constant $(2\pi)^{-1/2}$ is enclosed between $39894/10^5$ and $39895/10^5$, using $333/106<\pi<355/113$. For $x\ge0$, $\Phi(x)$ is enclosed by adding $1/2$ to these density bounds times
\[
 \sum_{k=0}^m\frac{(-1)^k x^{2k+1}}{2^k k!(2k+1)},
 \qquad m=101\ \text{(lower)},\quad m=100\ \text{(upper)},
\]
bounds that follow by integrating the alternating Taylor remainders of $e^{-u}$. The factor $e^{-h^2/2}$ in $r(h)$ is enclosed by Taylor polynomials of degrees 29 and 30. In the remaining exponential $e^v$, with $j=\max(1,\lceil |v|\rceil)$, the degree-17 Taylor lower bound for $e^{v/j}$ is raised to the power $j$. All rounding is outward, and the final comparison is of integers against $10^{10}$. The block strengths $u$ and the tangent point $h$ are exact rational inputs; they are close to the values that optimise these bounds, so the margin of the certificate over $10^5$ is about $10\%$.

\section{Companion programs}\label{app:code}

The programs are written in Python. Each certificate or check exits with an error if any check fails, including under \texttt{python -O}.
\begin{itemize}[leftmargin=*]
\item \texttt{certificate.py}: exact rational certificates for Proposition~\ref{prop:78} and the 319-candidate witness, the triangle decomposition and regularity of the 39-vertex blocks, and the automorphism counts of Corollary~\ref{cor:classes}.
\item \texttt{example78.py}: writes both 78-candidate tournaments as a labelled table of the 3003 pair winners and recounts their outdegrees and all 76,076 triples.
\item \texttt{check\_compression.py}: exact enumeration on five vertices (Section~\ref{sec:compression}).
\item \texttt{check\_cycles.py}: exhaustive verification of the fair-sign symmetry and exact calibration on every score class for $n\le6$.
\item \texttt{score\_test.py}, \texttt{kernel\_test.py}, \texttt{probit\_test.py}: the scan, spectral and targeted tests; \texttt{check\_scan.py} and \texttt{check\_spectral.py} check them against direct evaluation and on positive controls.
\item \texttt{power\_study.py}: the seeded simulation of Table~\ref{tab:power}.
\item \texttt{small\_reversals.py}: the Monte Carlo computations of Section~\ref{sec:small}.
\end{itemize}

\begin{thebibliography}{99}
\bibitem{harrison} M. Harrison-Trainor (2022).
An analysis of random elections with large numbers of voters.
\emph{Mathematical Social Sciences} \textbf{116}, 68--84.
\url{https://doi.org/10.1016/j.mathsocsci.2022.01.002}.
arXiv version: \url{https://arxiv.org/abs/2009.02979}.

\bibitem{close} J. H\k{a}z\l{}a, E. Mossel, N. Ross and G. Zheng (2020).
The probability of intransitivity in dice and close elections.
\emph{Probability Theory and Related Fields} \textbf{178}, 951--1009.
\url{https://doi.org/10.1007/s00440-020-00994-7}.

\bibitem{Mossel2010} E. Mossel (2010).
Gaussian bounds for noise correlation of functions.
\emph{Geometric and Functional Analysis} \textbf{19}, 1713--1756.
\url{https://doi.org/10.1007/s00039-010-0047-x}.

\bibitem{varin} C. Varin and D. Firth (2024).
Tractable ridge regression for paired comparisons.
\emph{Stat} \textbf{13}, e70028.
\url{https://doi.org/10.1002/sta4.70028}.

\bibitem{sst} N. Shah, S. Balakrishnan, A. Guntuboyina and M. Wainwright (2016).
Stochastically transitive models for pairwise comparisons: statistical and computational issues.
\emph{Proceedings of Machine Learning Research} \textbf{48}, 11--20.
\url{https://proceedings.mlr.press/v48/shahb16.html}.

\bibitem{saile} C. Saile and W. Suksompong (2020).
Robust bounds on choosing from large tournaments.
\emph{Social Choice and Welfare} \textbf{54}, 87--110.
\url{https://doi.org/10.1007/s00355-019-01213-6}.

\bibitem{SD2026} R. Singh and O. Davidov (2026).
Testing for lack of fit in paired comparison data.
\url{https://arxiv.org/abs/2604.00426}.

\bibitem{Lauritzen2008} S. L. Lauritzen (2008).
Exchangeable Rasch matrices. \emph{Rendiconti di Matematica e delle sue Applicazioni, Serie VII}
\textbf{28}, 83--95.

\bibitem{Iyer2016} T. Iyer (2016).
\emph{Random tournaments}. Bachelor of Philosophy (Honours) thesis,
Australian National University.
\url{https://tqft.net/web/research/students/TejasIyer/thesis.pdf}.

\bibitem{GT2016} A. Georgakopoulos and K. Tyros (2016).
The Bradley--Terry condition is $L_1$-testable.
\url{https://arxiv.org/abs/1609.05194}.

\bibitem{Thornblad2016} E. Th\"ornblad (2016).
Tournament limits: degree distributions, score functions and self-converseness.
\url{https://arxiv.org/abs/1611.09579}.

\bibitem{usvt} S. Chatterjee (2015).
Matrix estimation by universal singular value thresholding.
\emph{Annals of Statistics} \textbf{43}, 177--214.
\url{https://doi.org/10.1214/14-AOS1272}.

\bibitem{MS2025} A. Makur and J. Singh (2025).
Minimax hypothesis testing for the Bradley--Terry--Luce model.
\emph{IEEE Transactions on Information Theory} \textbf{71}, 9163--9202.
\url{https://doi.org/10.1109/TIT.2025.3605846}.

\bibitem{RBSS2022} C. Rastogi, S. Balakrishnan, N. B. Shah and A. Singh (2022).
Two-sample testing on ranked preference data and the role of modeling assumptions.
\emph{Journal of Machine Learning Research} \textbf{23}(225), 1--48.
\url{https://jmlr.org/papers/v23/20-1304.html}.
\end{thebibliography}
\end{document}
